\documentclass[twocolumn, 12pt]{article}
\usepackage[utf8]{inputenc}
\usepackage[spanish]{babel}
\usepackage{graphicx}
\usepackage{float}
\usepackage{amsmath}
\usepackage{amsfonts}
\usepackage{geometry}
\geometry{a4paper, margin=2.5cm}

\title{\Large{Simulación y Análisis Teórico de la Interacción Laser-Electrón en Regímenes de QED de Campo Fuerte Mediante Enfoque Bidipolar Optimizado}}

\author{Martina García, Juan Hernandez, Nathaly Restrepo\\ \textit{Introducción a la Física Subatómica, Departamento de Física} \\ \textit{Universidad Nacional de Colombia, Sede Bogotá}}
\date{\today}

\begin{document}

\twocolumn[
  \maketitle
  \begin{center}
  \begin{minipage}{0.85\textwidth}
  \small
  \textbf{Resumen} \\
  \vspace{0.2cm}
En este trabajo se presenta el marco teórico, el diseño experimental propositivo y la arquitectura computacional para simular la colisión de haces de electrones relativistas con pulsos láser ultra-intensos en el régimen de Electrodinámica Cuántica de Campo Fuerte (SFQED). Inspirados en el estudio de Olofsson y Gonoskov (Phys. Rev. A 106, 063512), nos enfocamos en maximizar el parámetro de no linealidad cuántica $\chi$ mediante una configuración de enfoque geométrico denominada \textit{onda bidipolar}, superando los límites convencionales de focalización parabólica $f/2$. Se formula el esquema metodológico para distinguir la señal de fotones emitidos a alto $\chi$ frente al fondo estocástico de emisiones a bajo campo e interacción multienfoque. Se estructuran las bases del código numérico en Python empleando propagación paraxial escalar mediante \textsf{LightPipes} (espectro angular y apodización) e integración vectorial focal (a ser ejecutado en el módulo computacional) y se dejan configuradas las secciones analíticas para la posterior incorporación de los datos de simulación.
\vspace{0.2cm}
\\
\textbf{Palabras clave:} SFQED, Parámetro de no linealidad cuántica $\chi$, Enfoque bidipolar, Reacción de radiación, Simulación PIC/Monte Carlo, LightPipes.

  \end{minipage}
  \end{center}
  \vspace{0.3cm}
]

\section{Introducción}

El desarrollo de instalaciones láser de potencia multipetavatio (PW) combinado con aceleradores de electrones lineales o de aceleración por campo de estela en plasma (LWFA) ha abierto una nueva frontera en la física subatómica: la Electrodinámica Cuántica de Campo Fuerte (SFQED, por sus siglas en inglés). En estos regímenes, el campo electromagnético experimentado por el electrón en su sistema de reposo se aproxima o supera el campo crítico de Schwinger:
\begin{equation}
E_{cr} = \frac{m^2 c^3}{e \hbar} \approx 1.32 \times 10^{18} \text{ V/m}
\end{equation}

La interacción de cargas en tales campos extremos está gobernada por el parámetro estocástico no lineal cuántico $\chi$, definido invariantemente como:
\begin{equation}
\chi = \frac{\gamma}{E_{cr}} \sqrt{\left( \mathbf{E} + \frac{\mathbf{v}}{c} \times \mathbf{B} \right)^2 - \left( \frac{\mathbf{E} \cdot \mathbf{v}}{c} \right)^2}
\end{equation}
notación equivalente al invariantismo tensor relativista $\chi = \frac{e\hbar}{m^3 c^4} \sqrt{-(F_{\mu\nu} p^\nu)^2}$.

Al alcanzar valores de $\chi \gtrsim 1$, los procesos de radiación dejan de ser descritos por la electrodinámica clásica (fórmula de Larmor-Liénard-Wiechert) y requieren la incorporación de la corrección estocástica por emisión de fotones de alta energía y la creación de pares electrón-positrón ($\text{e}^- \text{e}^+$) por el proceso de Breit-Wheeler no lineal. Para $\chi \ge 1600$, se teoriza la conjetura de Ritus-Narozhny sobre el colapso de la QED perturbativa.

No obstante, una dificultad experimental crítica radica en la deconvolución de la señal: debido a las variaciones tiro a tiro (\textit{shot-to-shot variations}) del parámetro de impacto entre el haz de electrones y el foco láser, así como a las pérdidas continuas de energía antes de alcanzar el punto de máxima intensidad, la señal proveniente de eventos de alto $\chi$ queda oscurecida por eventos dominantes de baja transferencia de momento.

El presente documento establece el marco teórico y la metodología analítica para reproducir y simular el experimento propuesto por Olofsson \& Gonoskov (2022), utilizando una focalización de onda bidipolar para maximizar $\chi$ hasta valores de $\chi \approx 5.25 \left[\frac{\epsilon}{1 \text{ GeV}}\right] \left[\frac{P}{1 \text{ PW}}\right]^{1/2} \left[\frac{1 \ \mu\text{m}}{\lambda}\right]$.

\section{Marco Teórico}

\subsection{Parámetro No Lineal Cuántico y Reacción de Radiación}
Cuando un electrón de alta energía ($\gamma \gg 1$) colisiona frontalmente con una onda electromagnética intensa, la aceleración transversal inducida produce radiación gamma intensa. El régimen de reacción de radiación cuántica modifica la trayectoria media y genera ensanchamiento estocástico del espectro energético de los electrones. La probabilidad de emisión por unidad de tiempo depende directamente de las funciones de Bessel modificadas $K_{2/3}$ a través de la tasa de dispersión Compton no lineal:
\begin{equation}
\begin{split}
W_{\text{rad}}(\gamma, \chi) = & \frac{\alpha m c^2}{\sqrt{3} \pi \hbar \gamma} \int_0^\infty \frac{5x^2 + 12x + 12}{(1+x)^3} \\
& \times K_{2/3}\left(\frac{2x}{3\chi}\right) dx
\end{split}
\end{equation}
donde $x = \frac{\hbar \omega}{m c^2 \gamma - \hbar \omega}$ representa la transferencia relativa de energía al fotón.

\subsection{Geometría de Enfoque Bidipolar}
Para maximizar el campo en el foco dada una potencia láser total $P$, la descomposición multipolar demuestra que la superposición sincrónica de ondas dipolares eléctricas y magnéticas ofrece la máxima densidad de energía puntual (Bassett, 1986).

Definiendo la onda bidipolar como la combinación de un dipolo eléctrico focalizado a lo largo del eje $z$ y un dipolo magnético perpendicular a lo largo del eje $x$ ($a=1/2, b=0$), se obtiene la amplitud máxima de campo relativista $a_{\text{max}} = e E / (m c \omega_0)$:
\begin{equation}
a_{\text{max}} \approx 550 \left[ \frac{P}{1 \text{ PW}} \right]^{1/2}
\end{equation}
lo que otorga una constante de proporcionalidad $\kappa \approx 5.25$ para el cálculo del $\chi_{\text{max}}$ alcanzable. Esta distribución óptima de intensidad antes de ser reflejada sobre el espejo parabólico de focalización sigue el perfil teórico de Sheppard-Larkin:
\begin{equation}
I^p(R) \propto \left[ \left(\frac{R}{2L}\right)^2 + 1 \right]^{-4}
\label{eq:sheppard}
\end{equation}
donde $L$ representa la distancia focal y $R$ la distancia radial al eje del sistema.

\subsection{Criterio de Selección de Señal por Ángulo de Desviación}
Para distinguir los fotones emitidos en la región de alto $\chi$ frente al fondo de dispersión de baja energía, se emplea el correlativo cinemático entre el ángulo de emisión $\alpha$ (ángulo formado entre la dirección del fotón emitido y el momento inicial del electrón) y la amplitud del campo local. Los fotones emitidos en el campo máximo satisfacen:
\begin{equation}
\alpha \sim \frac{a_{\text{max}}}{\gamma}
\end{equation}
Filtrando los eventos que cumplen simultáneamente $\hbar \omega > \frac{1}{2} m c^2 \gamma$ y $\alpha > 0.6 \frac{a_{\text{max}}}{\gamma}$, es posible aislar con alta fidelidad las interacciones SFQED del régimen supercrítico.

\section{Metodología Experimental y Diseño de la Simulación}

La implementación del estudio se divide en dos fases principales: el diseño conceptual del ensamble físico (Figura de colisionador láser-electrón) y la arquitectura de simulación computacional (en desarrollo por los demás miembros del equipo).

\subsection{Esquema del Experimento Propositivo}
El arreglo experimental simulado consta de:
\begin{enumerate}
    \item \textbf{Fuente de electrones:} Haz de electrones relativistas ($\epsilon = 1 - 10 \text{ GeV}$) producido por aceleración LWFA o acelerador convencional lineal, con una carga típica de $100 \text{ pC}$ y densidad $n_e \sim 10^{19} \text{ cm}^{-3}$.
    \item \textbf{Sistema óptico de focalización:} Espejo parabólico de apertura ultra-ancha ($f$-number $\sim 0.25$) acoplado a un espejo plano con perforación central para el paso del haz de electrones y salida de fotones gamma.
    \item \textbf{Espectrómetro y Detectores:} Sistema de detección angular de fotones de alta energía colocado en la dirección de propagación anterior para registrar distribuciones $d^2 N / d\alpha d\omega$.
\end{enumerate}

\subsection{Arquitectura del Modelo Computacional}
La simulación programada por el grupo de cómputo empleará un esquema de conformación óptica paraxial mediante la librería \textsf{LightPipes} e integración vectorial combinada con algoritmos Monte Carlo / PIC para los eventos estocásticos de SFQED:
\begin{itemize}
    \item \textbf{Conformación y Propagación del Haz:} Se modela la fuente láser del laboratorio (Ti:sapphire, $170\text{ fs}$, $13\text{ kW}$, $\lambda \in [390, 1040]\text{ nm}$) como un haz gaussiano con cintura $w_0$ y factor de calidad $M^2$. La propagación a través de lentes delgadas $t(r) = \exp[-i k r^2/(2f)]$ y apodizadores se efectúa mediante el algoritmo de espectro angular (FFT):
    \begin{equation}
    \begin{split}
    U(x,y,z+\Delta z) = & \mathcal{F}^{-1} \Big\{ \mathcal{F}\{U\} \\
    & \times \exp\left[ -i \Delta z \frac{k_x^2 + k_y^2}{2k} \right] \Big\}
    \end{split}
    \end{equation}
    
    \item \textbf{Integración de Trayectorias y Colisión Estadística:} Algoritmo Boris/Vay para la dinámica de Lorentz del electrón cruzando el foco con un parámetro de impacto $b$ muestreado de una distribución uniforme en un disco: $p(b) = 2b / b_{\text{max}}^2$. Se calcula el campo instantáneo $\chi(b,t) = 2\gamma |E(b,ct)| / E_{cr}$ y la fracción del tiempo en alto campo:
    \begin{equation}
    F = \frac{1}{T} \int \Theta\left[ \chi - \frac{\chi_{\text{max}}}{2} \right] dt
    \end{equation}
    
    \item \textbf{Métricas de Validación:} Comparación entre el perfil simulado $I_{\text{sim}}$ y el teórico de Sheppard-Larkin [Eq. (\ref{eq:sheppard})] mediante la correlación de Pearson $r$ y el Error Cuadrático Medio:
    \begin{equation}
    \text{RMSE} = \sqrt{ \frac{1}{N} \sum_{i=1}^N \left( I_{\text{sim}}(R_i) - I^p(R_i) \right)^2 }
    \end{equation}
    En la escala real del Laboratorio de Óptica Avanzada de la UNAL, el parámetro pico se estima en $\chi_{\text{max}} \approx 2 \times 10^{-5}$ por GeV.
\end{itemize}

\section{Resultados y Análisis}
\textit{(Esta sección queda reservada para la incorporación de las gráficas, distribuciones angulares $d^2 f / d\alpha d\chi$, espectros energéticos de fotones, curvas de $\kappa$ en función del $f$-number y tablas de sección eficaz efectiva $\sigma$ y razón de señal $\eta$ que serán computadas e ingresadas por el equipo de simulación.)}

\subsection{Distribución Angular y Energética de Fotones}
\textit{[Espacio reservado para Gráfica 1: Histograma tridimensional de eventos detectados en función de $\alpha / (a_{\text{max}}/\gamma)$ y energía del fotón].}

\subsection{Escalamiento del Parámetro $\chi_{\text{max}}$ y Eficiencia de Selección}
\textit{[Espacio reservado para Tabla 1: Comparación de la señal obtenida para polarización lineal vs. circular a diferentes potencias de láser $P = 1 \text{ TW} - 1 \text{ PW}$ y métricas de ajuste RMSE/$r$].}

\section{Conclusiones}
\textit{(Sección reservada para resumir los hallazgos cuantitativos derivados de los datos de la simulación. Se evaluará el cumplimiento de la detección a nivel de confianza de $3\sigma$ con el número de disparos calculado, la precisión del modelo óptico en LightPipes frente al campo bidipolar vectorial y la factibilidad del uso del enfoque bidipolar en experimentos de física subatómica de frontera).}

\begin{thebibliography}{99}

\bibitem{olofsson2022}
C. Olofsson and A. Gonoskov, \textit{Attaining a strong-field QED signal at laser-electron colliders with optimized focusing}, Phys. Rev. A \textbf{106}, 063512 (2022).

\bibitem{dipiazza2012}
A. Di Piazza, C. Müller, K. Z. Hatsagortsyan, and C. H. Keitel, \textit{Extremely high-intensity laser interactions with fundamental quantum systems}, Rev. Mod. Phys. \textbf{84}, 1177 (2012).

\bibitem{gonoskov2022}
A. Gonoskov, T. G. Blackburn, M. Marklund, and S. S. Bulanov, \textit{Charged particle motion and radiation in strong electromagnetic fields}, Rev. Mod. Phys. \textbf{94}, 045001 (2022).

\bibitem{bassett1986}
I. M. Bassett, \textit{Limit to concentration by focusing}, Opt. Acta \textbf{33}, 279 (1986).

\bibitem{sheppard1994}
C. Sheppard and K. Larkin, \textit{Optimal concentration of electromagnetic radiation}, J. Mod. Opt. \textbf{41}, 1495 (1994).

\bibitem{bula1996}
C. Bula \textit{et al.}, \textit{Observation of Nonlinear Effects in Compton Scattering}, Phys. Rev. Lett. \textbf{76}, 3116 (1996).

\bibitem{fedotov2017}
A. Fedotov, \textit{Conjecture of perturbative QED breakdown at $\alpha\chi^{2/3}\ge 1$}, J. Phys. Conf. Ser. \textbf{826}, 012027 (2017).

\bibitem{luxe2023}
LUXE Collaboration, \textit{Technical Design Report for the LUXE Experiment}, arXiv:2308.00515 (2023).

\end{thebibliography}

\end{document}